\documentclass{article}
\usepackage[UTF8]{ctex}
\begin{document}

\title{无理数}
\author{}
\date{}
\maketitle

\textbf{定义：} 不能表示为两整数之比的实数称为无理数。

\textbf{经典证明：} $\sqrt{2}$ 是无理数。
假设 $\sqrt{2}=p/q$（既约分数），则 $p^2=2q^2$。
故 $p$ 为偶数，设 $p=2k$，代入得 $4k^2=2q^2$，即 $q^2=2k^2$。
故 $q$ 亦为偶数，与既约矛盾。$\square$

\textbf{常见无理数：} $\pi$、$e$、$\sqrt{3}$、$\varphi$（黄金比例）等。

无理数的小数表示无限不循环，与有理数共同构成实数系。

\end{document}
